%% Priyadarshi Roy — Master Resume (ATS Optimised, 2 pages)
%% Compile with: pdflatex priyadarshi_roy_resume.tex (run twice)
%%
%% ── MASTER RESUME TAILORING GUIDE ────────────────────────────────────────────
%% AI / GenAI Engineer     : prioritise RAG, Study Advisor, Voice Assistant, GPT Researcher.
%% Computer Vision Engineer: prioritise Thesis, HandGX, Vision Transformers, Image Retrieval.
%% NLP Engineer            : prioritise Summarisation, BERT, Transformers, Semantic Search.
%% General ML Engineer     : prioritise Thesis, Study Advisor, Emotion Classification.
%%
%% LENGTH CONTROL — this file is tuned to exactly 2 pages. If you add content, cut in
%% this order: (1) "Earlier Projects", (2) the two weakest certifications,
%% (3) the third bullet of each project.
%%
%% ATS RULES OBSERVED — keep these if you edit:
%%   * standard section headings, all caps, one per line
%%   * no tables, no multi-column layout, no text boxes, no images
%%   * plain bullet lists only; dates right-aligned as plain text
%%   * consistent date format: "Mon YYYY -- Mon YYYY" or "YYYY -- YYYY"
%%   * spell acronyms out once, then abbreviate: e.g. Retrieval-Augmented Generation (RAG)
%%   * consistent British spelling throughout
%% ─────────────────────────────────────────────────────────────────────────────
\documentclass[10pt,a4paper]{article}
% ── Packages ──────────────────────────────────────────────────────────────────
\usepackage[a4paper,top=1.25cm,bottom=1.15cm,left=1.5cm,right=1.5cm]{geometry}
\usepackage[T1]{fontenc}
\usepackage[utf8]{inputenc}
\usepackage{lmodern}      % scalable outline fonts (clean PDF text extraction for ATS)
\usepackage{xcolor}
\usepackage{enumitem}
\usepackage{microtype}
\usepackage{hyperref}
% ── Styling ───────────────────────────────────────────────────────────────────
\definecolor{headerblue}{RGB}{31,119,180}
\definecolor{darkgray}{RGB}{60,60,60}
\hypersetup{
    colorlinks  = true,
    urlcolor    = headerblue,
    pdftitle    = {Priyadarshi Roy - Machine Learning Engineer},
    pdfauthor   = {Priyadarshi Roy},
    pdfsubject  = {Resume},
    pdfkeywords = {Machine Learning Engineer, NLP, LLM, RAG, Generative AI,
                   Computer Vision, Python, PyTorch, LangChain, Docker}
}
\pagestyle{empty}
\setlength{\parindent}{0pt}
\setlength{\parskip}{0pt}
% ── Section heading ───────────────────────────────────────────────────────────
\renewcommand{\section}[1]{%
  \par\vspace{7pt}%
  \noindent{\large\bfseries\color{headerblue}#1}\par\vspace{2pt}%
  \noindent\hrule\par\vspace{3pt}%
}
\setlist[itemize]{
    leftmargin = *,
    itemsep    = 1pt,
    topsep     = 2pt,
    parsep     = 0pt,
    label      = \textbullet,
}
% ── Macros ────────────────────────────────────────────────────────────────────
% \entry{Title}{Organisation / Location}{Date}
\newcommand{\entry}[3]{%
  \noindent\textbf{#1}\hfill\textit{\small #3}\\[0pt]%
  {\small\textit{\color{darkgray}#2}}\par\vspace{2pt}%
}
% \project{Name}{Technologies}{Date}
\newcommand{\project}[3]{%
  \noindent\textbf{#1}\hfill\textit{\small #3}\\[0pt]%
  {\small\textit{\color{darkgray}Technologies: #2}}\par\vspace{2pt}%
}
% \research{Title}{Institution}{Technologies}{Date}
\newcommand{\research}[4]{%
  \noindent\textbf{#1}\hfill\textit{\small #4}\\[0pt]%
  {\small\textit{\color{darkgray}#2}}\\[1pt]%
  {\small\textit{\color{darkgray}Technologies: #3}}\par\vspace{2pt}%
}
% \skill{Category}{Keywords} — plain text, no table, so parsers read it cleanly
\newcommand{\skill}[2]{%
  \noindent\textbf{#1:} #2\par\vspace{2pt}%
}
% \minorproject{Name}{Date}{Description}{Technologies}
\newcommand{\minorproject}[4]{%
  \item \textbf{#1} (#2): #3 \textit{#4}%
}
\newcommand{\vsm}{\vspace{5pt}}
% ══════════════════════════════════════════════════════════════════════════════
\begin{document}
% ── Header ────────────────────────────────────────────────────────────────────
\begin{center}
{\LARGE\bfseries\color{headerblue} PRIYADARSHI ROY}\\[3pt]
{\small\color{darkgray} Machine Learning Engineer \enspace|\enspace NLP \enspace|\enspace Generative AI \enspace|\enspace Computer Vision}\\[5pt]
\small
(+49) 1521-7847618 \enspace|\enspace
\href{mailto:proy9714@gmail.com}{proy9714@gmail.com} \enspace|\enspace
Cölbe, Germany \enspace|\enspace
\href{https://linkedin.com/in/priyadarshi-roy}{LinkedIn} \enspace|\enspace
\href{https://github.com/proy9714}{GitHub}
\end{center}
\vspace{1pt}
% ── Summary ───────────────────────────────────────────────────────────────────
\section{PROFESSIONAL SUMMARY}
Machine Learning Engineer with 3 years of industry experience designing and deploying AI
systems across Natural Language Processing, Generative AI, and Computer Vision. Experienced in
building Large Language Model (LLM) applications, Retrieval-Augmented Generation (RAG)
pipelines, semantic search systems, and production machine learning services using Python and
modern deep learning frameworks. Currently a \textbf{MSc Data Science} student at
\textbf{Philipps-Universität Marburg} (thesis due November 2026),
\textbf{available to start immediately}.
% ── Professional Experience ───────────────────────────────────────────────────
\section{PROFESSIONAL EXPERIENCE}
\entry{Machine Learning Engineer}
{Spring and River Pvt.\ Ltd., Kolkata, India}
{Aug 2021 -- Aug 2024}
\begin{itemize}
\item Developed production Natural Language Processing systems for automated document
understanding, including multi-document summarisation and highlight extraction pipelines that
replaced manual review of large document collections.
\item Built enterprise Retrieval-Augmented Generation (RAG) applications using GPT-3.5 Turbo,
embedding-based retrieval, LangChain, and vector search, enabling employees to query
organisational knowledge in natural language.
\item Designed modular LLM orchestration workflows and evaluation pipelines, improving the
reliability, maintainability, and scalability of Generative AI applications.
\item Developed and deployed Python REST APIs using Flask to serve machine learning models in
production, supporting AI-powered features across business applications.
\item Collaborated with engineering and product teams across the machine learning lifecycle,
translating business requirements into deployed AI solutions.
\end{itemize}
% ── Research Experience ───────────────────────────────────────────────────────
\section{RESEARCH EXPERIENCE}
\research
{Master's Thesis: Computer Vision for Historical Art Analysis}
{Philipps-Universität Marburg, Germany}
{PyTorch, Vision Transformers, Vision-Language Models, Human Pose Estimation, Image Retrieval}
{2026 -- Present}
\begin{itemize}
\item Researching Computer Vision foundation models for semantic understanding of human
figures, interactions, and visual structures in historical artworks.
\item Developing multimodal pipelines combining human pose estimation, object detection, image
embeddings, and geometric verification for structured analysis of paintings.
\item Investigating domain shift and representation learning when transferring vision models
from natural images to non-photographic artistic media, including scene graph generation and
knowledge-based retrieval.
\end{itemize}
% ── Projects ──────────────────────────────────────────────────────────────────
\section{PROJECTS}
\project{Study Advisor (RAG-Based Chatbot)}
{Python, FastAPI, LangChain, PostgreSQL, pgvector, Vue.js}
{2025}
\begin{itemize}
\item Architected a multi-stage Retrieval-Augmented Generation (RAG) pipeline for academic
advising, integrating intent classification, query decomposition, Text-to-SQL, semantic
retrieval, context summarisation, and LLM response generation.
\item Implemented hybrid retrieval in PostgreSQL using pgvector for semantic search and
pg\_trgm for fuzzy matching, removing the need for a separate vector database.
\item Benchmarked Llama 3.1-70B against Gemini 2.5-Flash with the RAGAS framework over 75
questions; Gemini led on faithfulness (0.65 vs.\ 0.44) and answer correctness (0.57 vs.\ 0.47),
informing the production model choice.
\end{itemize}
\vsm
\project{NLS Voice Assistant}
{Python, Faster-Whisper, Ollama, Docker Compose, LLMs, ASR, TTS}
{2026}
\begin{itemize}
\item Developed a task-oriented voice assistant for weather queries and multi-turn calendar
management, integrating speech recognition, local LLM orchestration, and text-to-speech into a
containerised Docker Compose pipeline.
\item Designed a scenario-based evaluation framework measuring Intent Accuracy, Slot Accuracy,
Temporal Resolution, and Task Completion Rate, achieving 100\% intent accuracy across all
evaluated scenarios.
\item Analysed failure modes including speech recognition errors and temporal reasoning
ambiguity, identifying targeted improvements for production reliability.
\end{itemize}
\vsm
\project{Multi-Label Emotion Classification}
{Python, PyTorch, Hugging Face Transformers, BERT, Scikit-Learn}
{2025}
\begin{itemize}
\item Developed a multi-label emotion classifier over BERT embeddings (768-dimensional) using a
Binary Relevance Logistic Regression architecture with balanced class weighting.
\item Applied MLSMOTE augmentation in embedding space to address severe label imbalance,
increasing training samples from 2,768 to 4,724 across 28 label combinations.
\item Achieved a weighted F1-score of 0.81, with best per-class performance on fear (F1 0.85)
and strong minority-class recall on anger (0.92).
\end{itemize}
\vsm
\noindent\textbf{Earlier Projects}\par\vspace{2pt}
\begin{itemize}
\minorproject{Distributed Discovery of Inclusion Dependencies}{2024}
{Fault-tolerant distributed system for discovering inclusion dependencies in datasets exceeding
single-node memory, built with the Akka Actor Framework and ported to Apache Spark.}
{Java, Akka, Apache Spark}
\minorproject{Comparative Study of ML Models for Text Classification}{2021}
{Evaluation of classical machine learning and deep learning approaches across multiple
algorithms and feature representations.}
{Python, TensorFlow, Keras, NLTK, Scikit-Learn}
\minorproject{HandGX -- ASL Gesture Recognition}{2019}
{CNN-based computer vision application for American Sign Language gesture recognition, covering
image preprocessing, feature extraction, and mobile deployment.}
{Python, OpenCV, TensorFlow, Keras, Kivy}
\end{itemize}
% ── Open Source Contributions ────────────────────────────────────────────────
\section{OPEN SOURCE CONTRIBUTIONS}
\entry
{GPT Researcher (GitHub Contributor)}
{\href{https://github.com/assafelovic/gpt-researcher}{github.com/assafelovic/gpt-researcher}}
{2024}
\begin{itemize}
\item Contributed feature implementations, bug fixes, and code improvements to GPT Researcher,
an open-source autonomous research agent built around LLM orchestration and web-based
information retrieval.
\item Worked with agentic AI workflows, Retrieval-Augmented Generation components, and
Python-based LLM tooling in a large open-source codebase, delivering changes through reviewed
pull requests.
\end{itemize}
% ── Technical Skills ──────────────────────────────────────────────────────────
\section{TECHNICAL SKILLS}
\skill{Programming}{Python, Java, C++, JavaScript, SQL}
\skill{Machine Learning}{PyTorch, TensorFlow, Keras, Scikit-Learn, Deep Learning,
Neural Networks, Model Evaluation, Feature Engineering, Machine Learning Lifecycle}
\skill{Generative AI}{Large Language Models (LLMs), Retrieval-Augmented Generation (RAG),
LangChain, Hugging Face Transformers, Prompt Engineering, Semantic Search, RAGAS Evaluation}
\skill{Computer Vision}{OpenCV, Vision Transformers, Human Pose Estimation, Object Detection,
Image Retrieval, Multimodal AI}
\skill{Backend and MLOps}{FastAPI, Flask, REST APIs, Docker, Docker Compose, Git, Linux, Ollama}
\skill{Data Engineering}{PostgreSQL, pgvector, MongoDB, MySQL, Apache Spark, Pandas,
Data Processing, Data Pipelines}
% ── Certifications ───────────────────────────────────────────────────────────
\section{CERTIFICATIONS}
\begin{itemize}
\item LangChain for LLM Application Development -- DeepLearning.AI
\item LangChain: Chat with Your Data -- DeepLearning.AI
\item Natural Language Processing with Classification and Vector Spaces -- Coursera
\item Build Basic Generative Adversarial Networks (GANs) -- Coursera
\item Python and Computer Vision Bootcamp -- ShapeAI
\end{itemize}
% ── Education ────────────────────────────────────────────────────────────────
\section{EDUCATION}
\entry
{Master of Data Science}
{Philipps-Universität Marburg, Germany}
{2024 -- Present}
\vspace{3pt}
\entry
{Master of Computer Science}
{St.\ Xavier's College (Autonomous), Kolkata, India}
{2019 -- 2021}
\vspace{3pt}
\entry
{Bachelor of Computer Science}
{St.\ Xavier's College (Autonomous), Kolkata, India}
{2016 -- 2019}
% ── Languages ────────────────────────────────────────────────────────────────
\section{LANGUAGES}
\begin{itemize}
\item English (C2)
\item Bengali (Native)
\item Hindi (C2)
\item Japanese (JLPT N3)
\item German (A2)
\end{itemize}
\end{document}